\documentclass[11pt]{article}

% Packages
\usepackage[utf8]{inputenc}
\usepackage{graphicx}
\usepackage{booktabs}
\usepackage{hyperref}
\usepackage{amsmath}
\usepackage{listings}
\usepackage{xcolor}

% Code listing style
\lstset{
  basicstyle=\ttfamily\small,
  breaklines=true,
  frame=single,
  numbers=left,
  numberstyle=\tiny\color{gray}
}

\title{Semantic Projection: Zero-Hydration UI for Agentic Systems}
\author{Aether UI Project}
\date{\today}

\begin{document}

\maketitle

\begin{abstract}
We present Aether UI, a zero-hydration interface framework that inverts the traditional client-side rendering paradigm. Instead of shipping JavaScript bundles (40-200KB) to hydrate virtual DOM trees, Aether projects pre-computed HTML slivers via Server-Sent Events to a minimal 2KB kernel. Our approach achieves sub-50ms time-to-interactive latency while maintaining design system portability and accessibility compliance. We demonstrate that semantic projection—pre-computed sliver graphs streamed via SSE—eliminates the hydration tax of generative UI frameworks while enabling edge-native deployment.
\end{abstract}

\section{Introduction}

Traditional UI frameworks (React, Vue, Angular) treat the browser as an application runtime. They ship JavaScript to the client, hydrate virtual DOM, reconcile state, and generate HTML at runtime. This paradigm introduces significant overhead:

\begin{itemize}
  \item \textbf{Bundle Size}: React + ReactDOM alone weighs 40-160KB gzipped
  \item \textbf{Hydration Tax}: Time spent parsing, compiling, and executing JavaScript
  \item \textbf{Reconciliation Cost}: Virtual DOM diffing on every state change
\end{itemize}

Aether UI inverts this paradigm. The browser becomes a projection surface, not an application runtime. The agent owns the interface DNA—structure, style, and logic live server-side. The browser receives pre-computed HTML slivers via SSE and renders them through a 2KB kernel.

\section{Architecture}

\subsection{The Three-Layer Model}

Aether employs a three-layer architecture:

\begin{enumerate}
  \item \textbf{Edge Layer}: Intent resolution, sliver lookup, SSE streaming
  \item \textbf{Adapter Layer}: Framework integration (Snap-In, Jump-In, Net-New)
  \item \textbf{Kernel Layer}: SSE client, CSS injection, slot hydration, focus preservation
\end{enumerate}

\subsection{The 4-Step Pipeline}

\begin{enumerate}
  \item \textbf{Ingestion}: Heavy JSON (100KB) → Canonical Model (1KB signal)
  \item \textbf{Intent Matching}: O(1) lookup: situation → sliver ID
  \item \textbf{DNA Hydration}: Template + Model → HTML (string replacement)
  \item \textbf{Delta Projection}: SSE stream → CSS variables + textContent
\end{enumerate}

\subsection{The 2KB Kernel}

The browser kernel has four responsibilities:

\begin{lstlisting}[language=JavaScript, caption=Kernel Core]
class AetherKernel {
  constructor(container, config) {
    this.connect();     // SSE connection
  }

  handlePulse(pulse) {
    this.applyVariables(pulse.vars);   // CSS injection
    this.injectContent(pulse.content); // Slot hydration
    this.restoreFocus();               // Focus preservation
  }
}
\end{lstlisting}

\section{Integration Modes}

\subsection{Snap-In (Coexistence)}
Host framework and Aether share the DOM. Aether claims a subtree while existing React/Vue/Angular continues.

\subsection{Jump-In (Replacement)}
Aether intercepts data, unmounts the existing framework, and claims territory. React props are transformed to Canonical Models.

\subsection{Net-New (Edge-Native)}
No host framework. Edge function pre-checks route eligibility, injects kernel, skips JavaScript bundle entirely.

\section{Benchmarks}

\begin{table}[h]
\centering
\begin{tabular}{lccc}
\toprule
Metric & React 18 & A2UI & Aether \\
\midrule
TTI & 200-500ms & 100-300ms & $<$50ms \\
JS Payload & 80-200KB & 15-40KB & 2KB \\
FCP & 150-400ms & 100-250ms & $<$30ms \\
Lighthouse & 70-85 & 75-90 & 95-100 \\
\bottomrule
\end{tabular}
\caption{Performance comparison across frameworks}
\end{table}

Our kernel achieves 819 bytes gzipped, well under the 2KB target. Time-to-interactive measurements show consistent sub-50ms latency for pure Aether applications.

\section{Design System Integration}

Aether doesn't ship HTML. Host systems provide semantic components. Agents possess them via CSS variables:

\begin{lstlisting}[language=HTML, caption=Slot-based possession]
<div class="eds-card" data-aether-possessable="true">
  <h3 data-aether-slot="title"></h3>
  <p data-aether-slot="body"></p>
</div>
\end{lstlisting}

Standard tokens (--tempo, --accent-color, --background-subtle) map across design systems via PSI (Presentation-Semantic Interface) mappings.

\section{Related Work}

\subsection{React Server Components}
RSC executes React on the server but still hydrates on the client. Aether eliminates hydration entirely.

\subsection{A2UI (Agent-to-UI)}
A2UI generates JSON specifications that adapters render. Aether projects pre-computed HTML, avoiding runtime generation latency.

\subsection{CopilotKit / Assistant UI}
These frameworks integrate AI into React applications. Aether inverts the model—agents own UI DNA, not the framework.

\section{Conclusion}

Aether UI demonstrates that semantic projection—pre-computed sliver graphs streamed via Server-Sent Events—eliminates the hydration tax of generative UI frameworks while maintaining design system portability, accessibility compliance, and edge-native latency.

The 2KB kernel, combined with SSE-based projection, achieves 10x improvement in time-to-interactive compared to React 18 and 40-100x reduction in JavaScript payload.

\section*{Availability}

Source code: \url{https://github.com/aether-ui/aether-ui}

\bibliographystyle{plain}
\bibliography{references}

\end{document}
